\section{An explicit two-state filtering benchmark}\label{sec:benchmark}
\begin{proposition}[A two-label stationary nonlinear filter]\label{prop:benchmark}
Consider a binary hidden state, transition matrix
$P=\left(\begin{smallmatrix}3/4&1/4\\1/4&3/4\end{smallmatrix}\right)$,
reports $N(1,1)$ and $N(-1,1)$ before that transition, and a protected binary Brier audit after it. Let $\Phi$ be the standard normal distribution function and define on $[1/2,3/4]$
\begin{equation}\label{eq:binary-fixed}
 G(q)=\frac14+\frac12\{q\Phi(1+\ell_q)+(1-q)\Phi(1-\ell_q)\},
 \qquad \ell_q=\tfrac12\log\frac q{1-q}.
\end{equation}
There is a unique solution $q_*=G(q_*)$. The homogeneous two-label observer uses programmed probabilities $q_*,1-q_*$, sends a report to the high label exactly when $y\ge-\frac12\log(q_s/(1-q_s))$, and outputs $(q_s,1-q_s)$. Its invariant label masses are $1/2,1/2$ and its exact stationary audit risk is
\begin{equation}\label{eq:binary-risk}
 R_\infty=2q_*(1-q_*).
\end{equation}
Every local common-grid observer for this experiment, with any number of labels, has a joint-ensemble minorization with $b=1$, $\gamma\ge1/120$. In particular, the displayed two-label observer has finite-average risk at most $2q_*(1-q_*)+480/n$ from every admissible initial probability and either initial label. The constants are deliberately conservative. This proposition does not assert global two-label optimality.
\end{proposition}
\begin{proof}
Write $p$ for the probability of hidden state one. The actual update is
\[
 F(y,p)=\tfrac14+\tfrac12\,\frac{pe^{2y}}{1-p+pe^{2y}},
 \qquad K=[1/4,3/4],\quad D=[1/8,7/8].
\]
The two cells are cut at $1/2$. A high report from programmed input $q$ is exactly $y\ge-\ell_q$. Reflection of both the hidden state and the report exchanges the two incoming labels and the two outgoing cells. With incoming masses $1/2,1/2$ and means $q,1-q$, the outgoing high mass is therefore $1/2$. Its unnormalized state-one contribution is the sum of the transition-prediction terms on the two high-report events. Conditional Gaussian tail integration yields its normalized probability
\[
 \tfrac14+\tfrac12\{q\Phi(1+\ell_q)+(1-q)\Phi(1-\ell_q)\}=G(q).
\]
Thus a fixed point of \eqref{eq:binary-fixed} is exactly the invariant calibration equation, not a pathwise approximate-posterior equation.

The identity $q\phi(1+\ell_q)=(1-q)\phi(1-\ell_q)$ cancels the derivative terms involving $\ell_q'$. Hence
\[
 G'(q)=\tfrac12\{\Phi(1+\ell_q)-\Phi(1-\ell_q)\}\in[0,1/2].
\]
Also $G(1/2)>1/2$, $G(3/4)<3/4$ and $G$ is increasing. It maps the closed interval into itself and is a strict contraction. This proves existence, uniqueness and geometric computation of $q_*$. Conditional Brier projection at its calibrated invariant ensemble gives \eqref{eq:binary-risk}. Certified iteration, rather than decimal approximation, defines the program constant; finite readout and comparison errors are treated in Section~\ref{sec:compilation}.

For the minorization, let $x\in[2/5,3/5]$ be an acquired output and put $w=2x-1/2\in[3/10,7/10]$. For any programmed $q\in K$, the inverse report is
$y=\frac12(\log(w/(1-w))-\log(q/(1-q)))$, so $|y|<1$. Its actual mixture density is at least $\phi(2)$, while $dy/dx=1/[w(1-w)]\ge4$. The acquired scalar density is therefore at least $4\phi(2)>1/8$ on this interval. With $\zeta$ uniform there, whose density is $5$, one may take $\epsilon=1/40$. For binary probabilities in $K$, the positive-cone comparison holds with $\kappa=1/3$. Lemma~\ref{lem:lift} gives $\gamma\ge1/120$ and $b=1$ without any erasure. Binary Brier loss is at most $2$, so \eqref{eq:entrance} costs at most $480/n$.

For reference, the acquired density is uniformly bounded on all of $K$, not only on the overlap interval. If $z=\log(w/(1-w))$, its actual density is a Gaussian mixture with variance $4$ and means bounded in absolute value by $B=\log7+2<4$, uniformly on $D$. Its density is at most $(2\sqrt{2\pi})^{-1}\exp[-((|z|-B)_+)^2/8]$. Since $dz/dx\le8e^{|z|}$, the acquired scalar density is at most $4e^{B+2}/\sqrt{2\pi}<1024$. This gives an explicit full-support density certificate; no Gaussian tail was removed. The induced Brier chart changes only the stated fixed norm constants.
\end{proof}
